\chapter{Introduction: one group, two channels, and what towers do to it}\label{chap:intro}

\section{The object}

Everything in this book starts from a finite $(q+1)$-regular graph $Y$ --- in the
arithmetic cases a quotient $\Gamma\backslash X$ of the Bruhat--Tits tree $X$ of
$\mathrm{PGL}_2$ over a non-archimedean local field with residue field of order $q$ --- and
from three integer matrices attached to it: the adjacency (Hecke) operator $T_p$, the
Laplacian $\Delta_p=(q+1)I-T_p$, and the non-backtracking (Hashimoto) matrix
$A^{\mathrm{nb}}$ on directed edges. The finite abelian group
\[
\Jac(Y)\;=\;\operatorname{Tor}\operatorname{coker}\Delta_p ,
\]
the \emph{sandpile} or \emph{critical} group, has order $\kappa(Y)$, the number of
spanning trees, and its structure --- not just its order --- is read off from the Smith
normal form of $\Delta_p$. On an arithmetic quotient, $\kappa(Y)=\frac1{|V|}\prod_f
\bigl(p+1-a_p(f)\bigr)$ is a product of inverse local $L$-values over the eigenforms
$f$ of level $Y$: Frobenius traces precipitate as torsion. This single identity
(Chapter~\ref{chap:I}, Theorem~1.3.5) is the seed of the programme; the word
``algorithmic'' in the title refers to the fact that the group, and everything built on
it, is computed exactly, by integer elimination, on every instance we discuss.

\section{The two channels}

There are two natural non-commutative algebras over $Y$. The \emph{boundary channel}
is the crossed product $C(\partial X)\rtimes\Gamma$, equivalently the Cuntz--Krieger
algebra of $A^{\mathrm{nb}}$; it carries a Satake-twisted KMS dynamics whose partition
function is exactly the local $L$-factor. The \emph{Laplacian channel} is the
$\Z[u^{\pm1}]$-module $\coker(1-uT_p+qu^2)$ whose Fitting ideal is the Ihara zeta
function and whose specialization at $u=1$ has torsion $\Jac(Y)$.

The first principle of the book is that \emph{these two channels do not see the same
thing}. The boundary algebra is arithmetically blind: its $K$-theory is
$\Z^r\oplus\Z/(r-1)$ with $r$ the rank of $\Gamma$, whatever the Hecke data
(Chapter~\ref{chap:I}, Theorem~1.3.7); and its KMS \emph{states} see only the Perron
branch of the spectrum, so tempered Satake parameters survive only as functionals
(Theorem~1.2.5, the weight-gap theorem 1.2.11). The Laplacian channel sees the
arithmetic but is a module, not an algebra. The obvious repair --- realize the Ihara
companion matrix $C=\bigl(\begin{smallmatrix}T_p&-qI\\ I&0\end{smallmatrix}\bigr)$ by a
non-negative matrix, hence by a Cuntz--Krieger algebra --- is impossible:
$\operatorname{Tr}C^2=-n(q-1)<0$, and shift equivalence preserves the nonzero spectrum,
so no dilation of any size works (Chapter~\ref{chap:II}). What \emph{can} be done is
to realize the group $\Z\oplus\Jac(Y)$ as $K_0$ of a Kirchberg algebra by a
$2n\times2n$ non-negative ``shadow'' matrix (Chapter~\ref{chap:VI}); this shares only
Bowen--Franks data, not spectrum, with the pencil, which is why it evades the
obstruction. In the projective limit the obstruction dissolves into a type
$\mathrm{III}_{1/q}$ factor with no trace, but so does $K_0$ (Chapter~\ref{chap:III}).

\section{Towers}

The second principle is that new content appears only along towers. For a finite
graph the Hecke--Dirac Fredholm module is linear algebra in $M_n(\Z)$; along a tower
$Y_0\leftarrow Y_1\leftarrow\cdots$ the groups $\Jac(Y_k)$ grow, and the growth law is
where the arithmetic lives.

\emph{Abelian towers} (Part~II, Chapters~\ref{chap:IV}, \ref{chap:IX}, \ref{chap:XIV}).
For a $\Z/\ell^k$ voltage tower the inverse limit of the Laplacian cokernels is
$\Lambda^n/D(1+T)\Lambda^n$, $\Lambda=\Z_\ell[[T]]$, with exact control and
characteristic ideal generated by $\Theta(T)=\det D(1+T)$; the Iwasawa invariants are
read off the Newton polygon. $\Theta$ has an exceptional zero of order exactly two at
the trivial character, and its second derivative is $-\kappa(Y_0)\|h_\alpha\|^2$, where
$h_\alpha$ is the harmonic representative of the voltage class: a graph
Mazur--Tate--Teitelbaum formula whose $L$-invariant $-\|h_\alpha\|^2$ is a rational
number. This $L$-invariant is, at once, the effective mass of the bottom Bloch band of
the $\Z$-cover, the diffusion constant of the accumulated voltage under simple random
walk, and the Albanese length of the class; it is \emph{not} a cyclic-cohomology index,
and Chapter~\ref{chap:XIV} proves that no such index exists on the commutative triple.

\emph{The Iwahori tower} (Chapters~\ref{chap:V}, \ref{chap:VII}, \ref{chap:XIII}). The
tower of iterated line digraphs (non-backtracking paths of length $k$) has no deck
group. Its sandpile orders obey $\operatorname{ord}_p|\Jac(Y_k)|=\mu p^k-k+\nu_0$ with
$\lambda=-1$, which no $\Z_p[[T]]$-module can produce; the $-k$ is an accumulated Euler
characteristic of a two-term degeneracy complex. The tower is governed instead by two
degeneracy maps $\tau$ (tail) and $\eta$ (head) whose incidence calculus
$\tau_*\eta^*=A_k$, $\eta^*\tau_*=A_{k+1}$, $\tau_*\tau^*=qI$ replaces the group ring.
The upper Hecke operator factors through the tail norm and therefore kills its kernel,
which makes the ``arc determinant'' exactly $q^{n_k(q-1)}$ and turns the growth law into
a snake-lemma identity. A universal (Cohn) localization of the degeneracy algebra carries
a $K_1$-class whose level evaluations are $\pm n_k\kappa(Y_k)$ and whose boundary is a
non-split extension of $\Jac(Y_k)$ by $\Z/n_k$; the analytic side of a ``main
conjecture'' for this class is the central open problem of the book.

\emph{Operator algebras of the tower} (Part~III). The tail map is an in-covering of
shadow graphs and induces unital embeddings $\mathcal O_k\to\mathcal O_{k+1}$ of the
level Kirchberg algebras, whose inductive limit is a Kirchberg algebra with
$K_1\cong\Z$ and $\operatorname{Tor}K_0$ the Pontryagin dual of the sandpile pro-module.
For $q=p$ prime the kernel of the tail norm on $\Jac(Y_{k+1})$ is exactly the
$p$-torsion, so the $p$-part of the pro-module is torsion-free of infinite rank --- a
third certificate that no commutative Iwasawa mechanism is at work. The head map is an
out-covering and induces no homomorphism, but a proper $C^*$-correspondence whose
$K_0$-index is $\eta_*$; the two Kasparov round trips are the Hecke operators.

\section{Rank two}

Part~IV replaces the tree by $\widetilde A_2$ buildings. The cubic Hecke pencil
$P(u)=I-uA_1+qu^2A_2-q^3u^3$ carries the arithmetic, but the $\Z/3$ type grading kills
the trace obstruction of rank one, and the realization problem becomes an instance of
the (open) weak Generalized Spectral Conjecture over $\Z$. Thick quotients have
$b_1=0$, so there are no abelian towers and no rank-two Iwasawa theory: the theory is
non-commutative from the start. The geodesic edge flow $L_E$ --- the rank-two
Hashimoto matrix --- carries the whole pencil through an explicit integer intertwiner,
and its complement is $\sqrt q$ times an orthogonal matrix, so the ``parahoric'' zeros
of the edge zeta lie on the critical circle for \emph{every} complex, Ramanujan or not;
the Ramanujan property lives entirely in the vertex pencil (Chapter~\ref{chap:XV}).

\section{Applications}

Part~V applies the sandpile/tower machinery outside number theory: the effective-mass
reading of the $L$-invariant becomes a statement about tree lattices in condensed
matter; the boundary-blindness theorem becomes a statement about what the boundary
algebra of $p$-adic AdS/CFT cannot see; the sign of the Iwasawa $\lambda$ certifies
whether a graph tower is generated by a normal subgroup; the trace obstruction gives
power-sum lower bounds on hidden eigenvalues in positive realizations; and, most
substantially, the critical group of a de Bruijn graph is shown to measure genome
assembly ambiguity beyond branch statistics, to split exactly along repeats, and --- for
two-copy repeats --- to be the critical group of a one-vertex map, with a real phage
genome as the running example.

\section{How to read the book}

Chapter~\ref{chap:I} is required for everything. Part~II needs Chapter~\ref{chap:II}
\S5 (the towers) and Chapter~\ref{chap:III} \S2 (the growth law). Part~III needs
Chapters~\ref{chap:V}--\ref{chap:VII} and \ref{chap:VI}. Part~IV needs
Chapter~\ref{chap:II} for the obstruction and Chapter~\ref{chap:VI} for the shadow
construction, and is otherwise self-contained. The applied chapters need only the
definitions of Chapter~\ref{chap:I}, \S3 and, for the biology, nothing beyond the
sandpile group of a digraph. The closing chapter collects every ``solved / open /
impossible'' verdict in one place, with the corrections that later chapters made to
earlier ones.
